\documentclass[10pt,landscape]{article}

\usepackage[margin=0.32in]{geometry}
\usepackage[T1]{fontenc}
\usepackage[utf8]{inputenc}
\usepackage{lmodern}
\usepackage{textcomp}
\usepackage{microtype}
\usepackage[table,dvipsnames]{xcolor}
\usepackage{array}
\usepackage{booktabs}
\usepackage{enumitem}
\usepackage{hyperref}
\usepackage{listings}
\usepackage{multicol}
\usepackage{tabularx}
\usepackage{titlesec}

\pagestyle{empty}
\setlength{\parindent}{0pt}
\setlength{\parskip}{1.5pt}
\setlength{\columnsep}{13pt}
\setlength{\columnseprule}{0.2pt}
\setcounter{secnumdepth}{0}

\definecolor{ToshBlue}{HTML}{0F4C81}
\definecolor{ToshInk}{HTML}{1F2937}
\definecolor{ToshMuted}{HTML}{5B6470}
\definecolor{ToshPanel}{HTML}{F4F7FB}
\definecolor{ToshRule}{HTML}{CBD5E1}
\definecolor{ToshCodeBg}{HTML}{F8FAFC}

\color{ToshInk}

\titleformat{\section}
  {\large\bfseries\color{ToshBlue}}
  {}
  {0pt}
  {}
\titlespacing*{\section}{0pt}{0.4em}{0.12em}

\setlist[itemize]{leftmargin=1.05em,itemsep=0.5pt,topsep=1.5pt,parsep=0pt,partopsep=0pt}

\lstdefinestyle{tosh}{
  basicstyle=\ttfamily\scriptsize,
  backgroundcolor=\color{ToshCodeBg},
  frame=single,
  rulecolor=\color{ToshRule},
  framerule=0.4pt,
  xleftmargin=0pt,
  xrightmargin=0pt,
  aboveskip=0.18em,
  belowskip=0.24em,
  columns=fullflexible,
  keepspaces=true,
  showstringspaces=false,
  breaklines=true,
  literate={°}{{\textdegree}}1
           {μ}{{\ensuremath{\mu}}}1
           {Ω}{{\ensuremath{\Omega}}}1
}
\lstset{style=tosh}

\newcommand{\sheettitle}[3]{
  {\LARGE\bfseries\color{ToshBlue} ToSh Cheat Sheet: #1\par}
  \vspace{0.08em}
  {\small\color{ToshMuted} #2 \hfill #3\par}
  \vspace{0.18em}
  {\color{ToshRule}\rule{\linewidth}{0.8pt}}
  \vspace{0.12em}
}

\newcommand{\note}[1]{
  \vspace{0.12em}
  \fcolorbox{ToshRule}{ToshPanel}{%
    \parbox{\dimexpr\linewidth-2\fboxsep-2\fboxrule\relax}{\footnotesize #1}%
  }
  \vspace{0.12em}
}

\newcommand{\tighttable}[2]{
  {\footnotesize
  \renewcommand{\arraystretch}{1.08}
  \begin{tabularx}{\linewidth}{@{}#1@{}}
  #2
  \end{tabularx}}
}


\begin{document}
\footnotesize
\sheettitle{Types, Modules, And Interop}{User-defined types, CLR access, reflection, and native interop basics}{README + examples + commands}

\begin{multicols*}{3}

\section{Using And Require}
\begin{lstlisting}
using System.IO
using System.IO = IO
\end{lstlisting}

\begin{lstlisting}
require "./utils.tosh"
require Inventory from "./toastlib.tosh"
require Inventory from "./toastlib.tosh" as Inv
require { Reporting, Utilities } from "./toastlib.tosh"
\end{lstlisting}

\begin{itemize}
\item \texttt{using} shortens CLR namespace references in the current lexical scope.
\item \texttt{require} loads a ToSh script/module once per session and imports names lexically.
\item Imported modules are accessed with dot notation, e.g. \texttt{Inv.sample\_items()}.
\end{itemize}

\section{Named Types}
\begin{lstlisting}
record Item(Name: string, Quantity: int)

enum StockState {
    Healthy, Low, Out
}
\end{lstlisting}

\begin{lstlisting}
class Animal {
    prop Name: string? = null
    prop Sound: string? = null

    Animal(name: string, sound: string) {
        $this.Name = $name
        $this.Sound = $sound
    }

    func Describe() -> string {
        return $"The {$this.Name} says {$this.Sound}"
    }
}
\end{lstlisting}

\note{Records are lightweight value objects; classes support constructors, methods, stored/computed properties, \texttt{static} members, and \texttt{shy} private members.}

\section{Modules}
\begin{lstlisting}
module Utilities {
    shy func hidden() { ... }

    export func banner(title: string) {
        writeline $"=== {$title} ==="
    }
}

Utilities.banner("Report")
\end{lstlisting}

\columnbreak

\section{CLR Interop}
\begin{lstlisting}
new System.Guid "d3b07384-d113-4ec6-a5d2-5b1d0d2e8c39"
new System.Random | call Next 1 100
call System.String Join ", " ["objects", "pipelines", "types"]
\end{lstlisting}

\begin{lstlisting}
"hello".ToUpper()
String.Join(" + ", ["a", "b"])
Math.Round(Math.PI, 4)
\end{lstlisting}

\note{Prefer fluent calls like \texttt{\$obj.Method()} and \texttt{TypeName.Method()}; use \texttt{call} or \texttt{call-method} when the invocation has to stay dynamic or pipeline-shaped.}

\tighttable{>{\ttfamily\raggedright\arraybackslash}p{0.27\linewidth}X}{
\toprule
Command & Use \\
\midrule
\texttt{new} & construct CLR or ToSh objects \\
\texttt{call} & invoke a CLR instance/static method \\
\texttt{call-method} & invoke a method by dynamic name \\
\texttt{cast} & convert values to a target type \\
\texttt{type-of} & return the runtime type \\
\bottomrule
}

\section{Reflection And Inspection}
\begin{lstlisting}
type-of $obj
describe-type list<int>
describe-type System.String
members System.DayOfWeek
methods $obj
constructors list<int>
\end{lstlisting}

\begin{lstlisting}
get-props $obj
get-prop $obj Name
set-prop $obj Name "Toast"
del-prop $obj Name
has-prop $obj Name
has-method $obj ToString
\end{lstlisting}

\begin{lstlisting}
inspect $obj
inspect --flat
name-of $obj
\end{lstlisting}

\note{Dynamic records and many runtime objects expose a shell-record interface, so reflection commands work well across both CLR objects and ToSh-native shapes.}

\section{Assembly Loading}
\begin{lstlisting}
load-assembly ./MyLib.dll
describe-type MyCompany.Tools.Widget
constructors MyCompany.Tools.Widget
\end{lstlisting}

\note{Load an assembly when a CLR type is not already visible to the session; after that, \texttt{new}, \texttt{cast}, and dot-syntax calls work normally.}

\columnbreak

\section{Pattern For Domain Code}
\begin{lstlisting}
module Inventory {
    record Item(Name: string, Quantity: int)
    enum StockState { Healthy, Low, Out }

    func state_of(item: Item) -> StockState {
        return match ($item.Quantity) {
            _ <= 0 => StockState.Out
            _ < 3 => StockState.Low
            default => StockState.Healthy
        }
    }
}
\end{lstlisting}

\begin{lstlisting}
var bread = new Inventory.Item("Bread", 2)
Inventory.state_of($bread)
\end{lstlisting}

\section{Native Interop Basics}
\begin{lstlisting}
bind native "./libcrypto.so" as Crypto {
    func sha256(data: ptr, len: ulong) -> ptr
}
\end{lstlisting}

Describe a C struct with \texttt{raw struct}. Padding is never
declared --- sequential layout aligns naturally, as C does.

\begin{lstlisting}
raw struct UtsName {
    sysname:  cstring[65]
    nodename: cstring[65]
    release:  cstring[65]
}

bind native "libc.so.6" as LibC {
    func uname(out buf: UtsName) -> ok
}

(LibC.uname()).release
\end{lstlisting}

\begin{lstlisting}
alloc buffer = 1024
write-buffer $buffer "hello"
read-buffer cstring $buffer
read-buffer long $buffer --at 8
size-of UtsName
offset-of UtsName release
forget $buffer
\end{lstlisting}

\tighttable{>{\ttfamily\raggedright\arraybackslash}p{0.29\linewidth}X}{
\toprule
Native tool & Purpose \\
\midrule
\texttt{raw struct} & declare a C memory layout \\
\texttt{bind native} & bind library exports \\
\texttt{alloc} / \texttt{forget} & manage unmanaged buffers \\
\texttt{read-buffer} / \texttt{write-buffer} & copy values to and from native memory \\
\texttt{size-of} / \texttt{offset-of} & inspect unmanaged layouts \\
\bottomrule
}

\note{The \texttt{native-*} names are canonical; the short forms above are blessed aliases. \texttt{native-free} has none --- \texttt{free} is the System memory command, so use \texttt{forget}, which unsets the variable and frees the buffer together.}

\note{\textbf{Use records for compact data, classes for behavior, CLR interop for existing .NET types, and native interop only when managed APIs are not enough.} ToSh lives directly on top of the CLR, so shell pipelines, user-defined types, and .NET APIs can coexist in the same script.}

\end{multicols*}
\newpage

\sheettitle{Types, Modules, And Interop}{Visibility, member forms, reflection workflow, and safer native patterns}{page 2 / front-back reference}

\begin{multicols*}{3}

\section{Visibility And Scope}
\begin{lstlisting}
module Utils {
    shy var counter = 0
    export func get-count() => $counter
    global var DEBUG = false
}
\end{lstlisting}

\tighttable{>{\ttfamily\raggedright\arraybackslash}p{0.27\linewidth}X}{
\toprule
Modifier & Effect \\
\midrule
\texttt{shy} & private to the current lexical scope or type \\
\texttt{export} & visible when a module is imported \\
\texttt{global} & session-wide binding outside local scope \\
\bottomrule
}

\section{Class Member Forms}
\begin{lstlisting}
class HexColor {
    prop Hex: string? {
        get => $this.hex_value
        set => $this.hex_value = $value
    }

    shy prop hex_value: string? = null
    static func Green() => new HexColor("#00FF00")
}
\end{lstlisting}

\begin{itemize}
\item Stored property: \texttt{prop Name: Type = value}
\item Computed property: \texttt{prop Magnitude: double => (...)}
\item Getter/setter property: use a block with \texttt{get} and \texttt{set}
\item Static members hang off the type: \texttt{HexColor.Green()}
\end{itemize}

\section{Reflection Workflow}
\begin{lstlisting}
types System.String
describe-type list<int>
members DateTimeOffset
methods "toast"
constructors System.Random
\end{lstlisting}

\section{Assembly Loading}
\begin{lstlisting}
load-assembly ./MyLib.dll
types Widget
describe-type MyCompany.Tools.Widget
constructors MyCompany.Tools.Widget
\end{lstlisting}

\columnbreak

\section{Object Mutation Tools}
\begin{lstlisting}
var original = {| Name = "Bread", Qty = 2 |}
var copy = (clone $original)
set-prop $copy Qty 5 | ignore
get-prop $copy Qty
del-prop $copy Qty
\end{lstlisting}

\begin{lstlisting}
has-prop $copy Name
has-method "hello" ToUpper
get-props $copy
get-methods "toast"
\end{lstlisting}

\note{Dynamic records are great for ad-hoc pipeline shapes; records and classes are better when you want named, reusable domain types.}

\section{Module Import Patterns}
\begin{lstlisting}
require "./utils.tosh"
require Inventory from "./toastlib.tosh"
require Inventory from "./toastlib.tosh" as Inv
require { Reporting, Utilities } from "./toastlib.tosh"
\end{lstlisting}

\begin{itemize}
\item \texttt{require} is lexical and cached once per session.
\item \texttt{using} only aliases CLR namespaces; it does not run scripts.
\item Use module qualification when you want APIs to stay discoverable in larger scripts.
\end{itemize}

\section{CLR Call Patterns}
\begin{lstlisting}
new System.Random().Next(1, 100)
"hello".Replace("h", "y")
String.Join(", ", ["a", "b", "c"])
Math.Round(Math.PI, 3)
\end{lstlisting}

\begin{lstlisting}
call System.String Join ", " ["x", "y"]
call-method "hello" Replace "llo" "y"
\end{lstlisting}

\section{Choose The Surface}
\tighttable{>{\ttfamily\raggedright\arraybackslash}p{0.25\linewidth}X}{
\toprule
Tool & Prefer it when \\
\midrule
\texttt{record} & you want immutable value-shaped data \\
\texttt{class} & you need behavior, methods, or private state \\
\texttt{module} & you are grouping functions and types into an API \\
CLR type & .NET already has the exact thing you need \\
native ABI & managed APIs are not enough \\
\bottomrule
}

\columnbreak

\section{Native Signature Patterns}
\begin{lstlisting}
bind native "libc.so.6" as LibC {
    func abs(int) -> int callconv cdecl
    func strlen(string) -> nuint
    func memcpy(nint, nint, nuint) -> nint
}
\end{lstlisting}

\tighttable{>{\ttfamily\raggedright\arraybackslash}p{0.29\linewidth}X}{
\toprule
Native type & Typical use \\
\midrule
\texttt{nint} / \texttt{nuint} & pointer-sized ints \\
\texttt{ptr} & raw buffer pointers \\
\texttt{cstring} & C string interop surfaces \\
\texttt{cstring[n]} & inline \texttt{char[n]} inside a raw struct \\
\texttt{buffer[n]} & output-string pair (pointer + length) \\
\texttt{out} / \texttt{ref} & by-reference native parameters \\
\bottomrule
}

\section{Return Conventions}
\begin{lstlisting}
func sysinfo(out i: SysInfo)      -> ok
func sysconf(name: int)           -> count
func gettimeofday(out t: TV, nint)-> auto
func getloadavg(out a: double[3],
                n: int) -> int where (_ == 3)
\end{lstlisting}

\tighttable{>{\ttfamily\raggedright\arraybackslash}p{0.24\linewidth}X}{
\toprule
Return & Success is \\
\midrule
\texttt{ok} & \texttt{0}; anything else throws \\
\texttt{count} & \texttt{>= 0}, and is the value \\
\texttt{auto} & always; projects out params \\
\texttt{where (\_)} & your predicate over \texttt{\_} \\
\bottomrule
}

A checked call yields its \texttt{out} parameter directly. Failure throws
\texttt{NativeError} with \texttt{Errno} and \texttt{ErrnoName}.

\section{Buffer Workflow}
\begin{lstlisting}
alloc src = 6
alloc dst = 6
write-buffer $src "toast"
LibC.memcpy($dst, $src, 6)
read-buffer cstring $dst
forget $src $dst
\end{lstlisting}

\begin{lstlisting}
alloc counter = long
write-buffer $counter 42 --at 0
read-buffer long $counter
forget $counter
\end{lstlisting}

\section{Layout Inspection}
\begin{lstlisting}
size-of int
size-of UtsName
offset-of SysInfo totalram
\end{lstlisting}

\section{Native Safety Rules}
\begin{itemize}
\item Prefer \texttt{cstring} or pointer types for native string surfaces.
\item Never declare padding --- sequential layout aligns as C does.
\item Use \texttt{byte}, not \texttt{bool}: \texttt{bool} marshals as a 4-byte Win32 \texttt{BOOL}.
\item \texttt{out} is engine-allocated; \texttt{ref} means the callee also reads.
\item \texttt{forget} every buffer you \texttt{alloc}.
\end{itemize}

\note{\textbf{Rule of thumb:} stay in CLR land unless you truly need a native ABI. When you do cross that boundary, declare the layout with \texttt{raw struct}, state the success convention, and let the engine own the memory.}

\end{multicols*}
\end{document}
